\section{Conclusion}
\label{sec:conclusion}

This study presents a fuller but still bounded method-paper evaluation of a JEPA-style
latent-surprise detector for gameplay glitch detection. On the frozen pair-disjoint TempGlitch
follow-up, the best recorded LeWM row remains stronger than the simple baselines and stronger than
the exact-support learned baselines in observed AUROC. At the same time, the strongest
learned-baseline comparison is statistically inconclusive, the negative support and calibration
support remain small, and the operating-point false-positive rate is high.

The broader evidence lanes sharpen the method story rather than simply enlarging it. K2 provides a
bounded public-benchmark slice on which the current LeWM rows underperform validated simple and
learned-baseline rows. K3 provides a controlled mechanistic readout with no observed SIGReg
prediction-loss benefit and no stable real-action-over-zero-action gain. Qualitative timelines
improve interpretability but do not unlock temporal localization because span labels are absent.

The resulting contribution is an artifact-backed, leakage-aware detector evaluation with explicit
positive and negative findings. It does not establish broad superiority, cross-game
generalization, temporal localization, deployment readiness, or locked-test performance. The next
submission step is operational rather than scientific: run the official LLNCS build in Overleaf
or the Springer author kit, record bibliography/reference warnings, verify anonymization and PDF
metadata, and complete external similarity screening.
